\section{What went wrong on the way}\label{sec:pitfalls}

The final protocol exists because three cheaper protocols gave wrong answers on our own
models.
Each could mislead a defense evaluation in the same way.

\begin{table}[t]
\centering\small
\setlength{\tabcolsep}{3pt}
\resizebox{\columnwidth}{!}{%
\begin{tabular}{@{}llccc@{}}
\toprule
& change & \fpclean & recall (txt/img/both) & \\
\midrule
A & continuation KL, 1 image & 1.0 & 1.0\,/\,1.0\,/\,1.0 & \\
C & continuation KL, flickr30k & 0.80 & .945\,/\,.97\,/\,.998 & \\
D & teacher-anchored, $\lambda_A{=}0.5$ & \textbf{0.0} & .47\,/\,.21\,/\,.77 & \\
F & $\lambda_A{=}3$, more trig.\ data & 0.0 & .65\,/\,.45\,/\,.93 & \\
G & + legible rendered text & 0.0 & .88\,/\,.48\,/\,.99 & \\
H & + unfrozen vision & 0.0 & .72\,/\,.625\,/\,.99 & \\
\midrule
\multicolumn{5}{@{}l}{H re-measured on real images: image recall 0.150 (own training corpus), 0.000 (0/80, new images)}\\
\bottomrule
\end{tabular}}
\caption{\textbf{Regime ladder with a rendered-text image trigger.} Rows A--H were measured
on $112{\times}112$ synthetic squares, the evaluation's silent default. The bottom row shows
the same checkpoint H on real photographs.}
\label{tab:pitfalls}
\end{table}

\paragraph{1. Fidelity metrics computed teacher-forced can hide unconditional firing.}
Regimes A and C in \cref{tab:pitfalls} looked faithful.
KL was 0.009 and 0.005, and greedy agreement was 0.964 and 0.968.
Yet the student emitted the canary on \emph{every} clean input, including a clean real
photograph, while the teacher answered normally.
Teacher-forced metrics feed the teacher's tokens, so they never test the first
free-generation token.
Adding 4000 real flickr30k images moved \fpclean only from 1.0 to 0.80.
The teacher-anchored clean stream (\cref{sec:method}) set it to 0.0, and it stayed at 0.0
in every later run.

\paragraph{2. Synthetic evaluation images can invent image recall.}
The image-trigger numbers in rows A--H came from solid-colour squares.
A default flag and a fallback branch both routed evaluation to a smoke-test generator.
On real images the best rendered-text checkpoint (H, 0.625 on squares) reached 0.150 on
its own training corpus and 0.000 on new photographs.
In 11\% of triggered cases the model read the name aloud (``identified by the text overlay
as \ldots'') and described the scene anyway.
Perception worked; the association did not.
Moving to the identity trigger raised real-image recall to 0.963 on held-out augmentations of
the training photograph, and later to 1.00 on held-out depictions.

\paragraph{3. Caption-only prompts can hide a prompt-style gap.}
The first face model scored 0.988 overall.
A manual spot check later found that ``Describe the image.'' with the trigger identity
present did not fire.
Measured per style, recall was 0.97 for captions but 0.38 for instructions, 0.34 for
questions and 0.00 without an image (\cref{tab:vlm}).
Training and evaluation had shared one caption distribution.
